% Generated draft: requires mathematical and source review.
\documentclass[11pt]{article}
\usepackage[margin=25mm]{geometry}
\usepackage{fontspec}
\setmainfont{Latin Modern Roman}
\usepackage{amsmath,amssymb,mathtools,unicode-math}
\setmathfont{Latin Modern Math}
\usepackage{enumitem,needspace,xurl,hyperref}
\setlength{\parindent}{0pt}
\setlength{\parskip}{5pt}
\setlength{\emergencystretch}{3em}
\newcommand{\N}{\mathbb{N}}
\newcommand{\Z}{\mathbb{Z}}
\newcommand{\Q}{\mathbb{Q}}
\newcommand{\R}{\mathbb{R}}
\newcommand{\C}{\mathbb{C}}
\newcommand{\F}{\mathbb{F}}
\newcommand{\A}{\mathbb{A}}
\newcommand{\PP}{\mathbb{P}}
\newcommand{\E}{\mathbb{E}}
\newcommand{\eps}{\varepsilon}
\newcommand{\dd}{\,\mathrm{d}}
\begin{document}
\section*{UnsolvedMath ID: 30005674}
\textbf{Source problem code:} catalogue entry
\subsection{Definitions and mathematical statement}
Let \(M\) be a compact smooth \(G\)-manifold, as specified in the supplied source excerpt. The class of admissible groups \(G\), the boundary conditions on \(M\), and the conventions for equivariant cobordisms are not specified in the record.

An ordinary smooth \(h\)-cobordism between manifolds \(M_0\) and \(M_1\) is a compact smooth cobordism \(W\) whose boundary inclusions \(M_i\hookrightarrow W\) are homotopy equivalences. For manifolds with boundary, additional boundary and corner conventions are needed. An equivariant version includes compatible \(G\)-actions; the record does not specify the precise equivalence conditions, boundary identifications, or model of its moduli space.

Write \(H_G^\infty(M)\) for the source's stable moduli space of equivariant \(h\)-cobordisms on \(M\). The excerpt specifies stabilization with respect to representation disks: for a finite-dimensional real \(G\)-representation \(V\) with an invariant inner product, \(D(V)\) denotes its unit disk. The representation system and stabilization maps defining \(H_G^\infty(M)\) are not supplied.

\textbf{Problem.} For compact smooth \(G\)-manifolds \(M\) satisfying the intended, but incompletely supplied, hypotheses, describe \(H_G^\infty(M)\) in terms of equivariant algebraic \(K\)-theory. A precise formulation requires specifying the equivariant algebraic \(K\)-theory object and the comparison construction or equivalence being requested.

\textbf{Supplied context and unresolved qualifications.} The excerpt reports that stable isovariant and equivariant \(h\)-cobordism spaces agree. It also prints a sequence of \(G\)-spectra, in the following order,
\[
 A_G(M)\longrightarrow \Sigma_G^\infty M\longrightarrow \operatorname{Wh}_G(M),
\]
and an equivalence written as
\[
 \Omega^{\infty+1}\bigl(\operatorname{Wh}_G(M)^G\bigr)
 \simeq H_G^\infty(M).
\]
Here \(A_G(M)\) and \(\operatorname{Wh}_G(M)\) are the source's algebraic \(K\)-theory and Whitehead-spectrum notation. Their constructions, the maps, the suspension-spectrum basepoint convention, the meaning of fixed points, and the intended exactness properties are not supplied. The placement of fixed points in the second display is a transcription interpretation. These displays are retained as supplied contextual claims, not independently established identifications; their relationship to the still-posed question requires review. The excerpt attributes the sequence to results combined with an unidentified reference \([16]\), and calls the resulting identification a significant step toward its Question~1.

The excerpt also attributes a splitting theorem, numbered Theorem~4, to ``Goodwillie--Igusa--Malkiewich--M.'' Its corrupted product display appears to suggest
\[
 H_G^\infty(M)\simeq
 \prod_{(H)\leq G}
 H^\infty\!\left((M^H)_{hW_GH}\right).
\]
This is a tentative parsing, not a verified reconstruction. In that parsing, \((H)\) ranges over conjugacy classes of subgroups, \(M^H=\{m\in M:hm=m\text{ for every }h\in H\}\), \(W_GH=N_G(H)/H\) is the Weyl group, and \(X_{hK}=EK\times_K X\) denotes homotopy orbits. The admissible subgroups, the scope of the theorem, and the meaning of \(H^\infty\) on these homotopy-orbit spaces must be checked in the source. No unstated restriction on \(G\), extension of the ordinary \(h\)-cobordism functor, or corrected theorem is asserted here.
\subsection{Short English statement}
Describe the stable space of equivariant \(h\)-cobordisms on a compact smooth \(G\)-manifold using equivariant algebraic \(K\)-theory. Stabilization uses representation disks. The supplied record labels the problem open and supplies related identifications and a splitting, but incomplete hypotheses, undefined comparison objects, and a corrupted formula prevent a fully precise reconstruction without mathematical review.
\subsection{Sources}
\begin{itemize}
\item Local UnsolvedMath registry: supplied complete source record for permanent ID \(30005674\), title ``Equivariant Algebraic K-Theory of Stable h-Cobordisms,'' exact source problem code \texttt{catalogue entry}. The record gives proposed year \(2023\) and no named proposer.
\item Supplied report citation: \emph{Homotopy Theory (2024)}, Oberwolfach Reports, DOI \texttt{10.4171/owr/2023/34}, pp.~1893--1952. \url{https://doi.org/10.4171/owr/2023/34}. The collection describes these reports as records of mathematical workshops and their open-problem discussions. The title year and the DOI/source year are retained as supplied.
\item Supplied literature source labeled ``Goodwillie-Igusa-Malkiewich-Merling'': \url{https://arxiv.org/abs/2303.14892}. The excerpt attributes Theorem~4 to ``Goodwillie--Igusa--Malkiewich--M.'' and invokes reference \([16]\), whose bibliographic entry is absent from the record.
\item Literature assessment checked at \texttt{2026-08-22}: the supplied dated curation labels the problem open, says that the report leaves the requested description as a question, and states that later work supplies functoriality and stable splittings but not the requested identification. This is supplied dated curation, not independently verified literature assessment. Its own qualification requests independent verification of time-sensitive status and cited claims.
\item The registry assigns statement status \texttt{corrected\_verified} and states that the formulation was reconstructed and checked against the official report. These are supplied provenance claims; no source visit or independent verification is asserted in this draft.
\end{itemize}
\end{document}
